\chapter{Appendix to Chapter 4: Secondary Derivations, Proofs, and Reproducible Computation}
\label{app:ch4_proofs}

This appendix separates technical detail from the intuitive presentation in
Chapter~\ref{chap:model2}. It provides proof sketches, the complete
computational procedure, and the stylized input data used to generate every
numerical result in the chapter.

\section{Phase I Derivations and Proofs}

\subsection{Feasibility of Stewardship-Network Operation}

Equation~\eqref{eq:con_retailer} requires total flow
$D(S)=\sum_{i\in I(S)}D_i(S)$ to pass through both the consolidation and
processor echelons. If opened-center capacity or eligible-processor capacity is
below $D(S)$, flow conservation makes the model infeasible. Conversely, under
complete connectivity, the required flow can be distributed across any opened
centers and eligible processors whose aggregate capacities satisfy
\eqref{eq:ch4_feasibility_condition}. This proves
Proposition~\ref{prop:ch4_feasibility}.

\subsection{Value of Additional Processor Access}

Let $e'$ weakly expand processor access relative to $e$. Every solution feasible
under $e$ remains feasible under $e'$ because no access constraint is tightened.
The minimum over the larger feasible set cannot exceed the minimum over the
smaller set. Therefore $c(S;e')\leq c(S;e)$, proving
Proposition~\ref{prop:ch4_access_value}. The inequality can be an equality when
the newly available processor is not selected.

\subsection{Average Cost on a Fixed-Capacity Interval}

\begin{proposition}[Average-Cost Behavior Between Capacity Thresholds]
\label{prop:ch4_average_cost_behavior}
Suppose coalition volume $Q$ changes over an interval on which the same
facilities and processors remain active, total fixed cost is $F>0$, and the
unit net variable cost is $\bar v$. Then
\begin{equation}
\label{eq:ch4_average_cost_behavior}
\mathrm{UC}(Q)=\bar v+\frac{F}{Q},
\qquad
\frac{\partial\mathrm{UC}(Q)}{\partial Q}=-\frac{F}{Q^2}<0.
\end{equation}
At a capacity or access threshold, a new binary decision may instead create a
discrete increase or decrease in average cost.
\end{proposition}

When the opened facilities and unit net cost $\bar v$ remain unchanged, total
cost is $F+\bar vQ$. Dividing by $Q$ and differentiating gives
$\mathrm{UC}(Q)=\bar v+F/Q$ and
$\partial\mathrm{UC}/\partial Q=-F/Q^2<0$. A new binary facility decision
changes $F$ and can create an upward jump. This proves
Proposition~\ref{prop:ch4_average_cost_behavior}.

\section{Phase II Derivations and Proofs}

\subsection{Operational Anatomy of a Marginal Contribution}

For a predecessor coalition $T\subseteq N\setminus\{x\}$, player $x$'s
marginal cost contribution is
\begin{equation}
\label{eq:ch4_marginal_cost}
\Delta_x c(T)=c(T\cup\{x\})-c(T).
\end{equation}
For interpretation, this change can be organized as
\begin{equation}
\label{eq:ch4_marginal_decomposition}
\begin{aligned}
\Delta_x c(T)=&
\Delta_x C^{flow}(T)
+\Delta_x C^{fixed}(T)
+\Delta_x C^{access}(T)\\
&+\Delta_x C^{gov}(T)
-\Delta_x V^{rec}(T).
\end{aligned}
\end{equation}
The five terms represent changes in flow cost, active fixed facilities,
processor access, platform administration, and recovered-material credit.
This is an accounting interpretation of two optimized coalition costs, not an
additional separability assumption imposed on the MILP.

\subsection{Shapley Efficiency}

For any ordering of the players, marginal contributions telescope from
$c(\emptyset)=0$ to $c(N)$. Averaging the telescoping sum over all player
orderings leaves the sum unchanged, so
$\sum_{x\in N}\phi_x(c)=c(N)$. This proves
Proposition~\ref{prop:ch4_efficiency}.

\subsection{Counterexample to Automatic Individual Rationality}

Consider a two-player cost game with
$c(\emptyset)=c(\{1\})=c(\{2\})=0$ and $c(\{1,2\})=2$. Symmetry and efficiency
give $\phi_1(c)=\phi_2(c)=1$, but $1>c(\{1\})=c(\{2\})=0$. Shapley efficiency
therefore does not imply individual rationality. This proves
Proposition~\ref{prop:ch4_ir_not_automatic}.

\subsection{Core Stability}

For a cost game, a proper coalition blocks allocation $a$ exactly when its
assigned payment exceeds what it can achieve alone:
$\sum_{x\in S}a_x>c(S)$. This is equivalent to
$\sigma(S;a)<0$. Thus an efficient allocation is unblocked if and only if all
proper-coalition slacks are nonnegative. This proves
Proposition~\ref{prop:ch4_core_depends_phase1}.

\subsection{Hybrid Efficiency}

\begin{proposition}[Hybrid Allocation Preserves Budget Balance]
\label{prop:ch4_hybrid_efficiency}
If $\sum_{x\in N}p_x=1$, then the hybrid allocation in
Equation~\eqref{eq:ch4_hybrid_allocation} satisfies
\begin{equation}
\label{eq:ch4_hybrid_efficiency}
\sum_{x\in N}\psi_x(\theta)=c(N)
\qquad \forall \theta\in[0,1].
\end{equation}
\end{proposition}

Summing Equation~\eqref{eq:ch4_hybrid_allocation} over players gives
\[
\sum_x\psi_x(\theta)
=\theta\sum_x\phi_x(c)+(1-\theta)c(N)\sum_xp_x.
\]
Shapley efficiency and $\sum_xp_x=1$ reduce the expression to $c(N)$. This
proves Proposition~\ref{prop:ch4_hybrid_efficiency}.

\subsection{Non-Smooth Allocation Changes}

\begin{proposition}[Sufficient Condition for a Non-Smooth Allocation Change]
\label{prop:ch4_nonlinear_allocation}
Suppose a facility-capacity or processor-access threshold changes the optimized
cost of at least one coalition by a nonzero discrete amount, while the weighted
marginal-cost changes in the Shapley sums do not cancel exactly for every
player. Then at least one player's Shapley allocation changes non-smoothly at
that threshold.
\end{proposition}

Each Shapley allocation is a weighted linear combination of coalition-cost
differences. If a threshold creates a nonzero discrete change in at least one
coalition cost, it changes the relevant marginal contributions. Unless those
weighted changes cancel for every player, at least one Shapley allocation has a
non-smooth change. This proves Proposition~\ref{prop:ch4_nonlinear_allocation}
under its stated non-cancellation condition.

\section{Structural and Comparative-Static Results}
\label{app:ch4_structural}

Throughout this section $a(S)=\sum_{x\in S}a_x$ for a vector $a\in\mathbb{R}^N$,
``proper coalition'' means $\emptyset\neq S\subsetneq N$, and $\epsilon^*$ is
the least-core value of \eqref{eq:ch4_least_core} without the restriction
$a_x\geq0$.

\subsection{Additive Cost Components}

If $y$ is budget balanced for $c$ and $y(S)\leq c(S)$ for every proper $S$,
then $y+a$ is budget balanced for $c'=c+a(\cdot)$ and
$(y+a)(S)=y(S)+a(S)\leq c'(S)$; the converse holds by subtracting $a$. Hence
the core of $c'$ is the core of $c$ shifted by $a$. The same shift maps every
feasible pair $(y,\epsilon)$ of the least-core program for $c$ to a feasible
pair $(y+a,\epsilon)$ for $c'$ with the same $\epsilon$, so
$\epsilon^*(c')=\epsilon^*(c)$; applied to program
\eqref{eq:ch4_cost_of_stability} it shows that the Cost of Stability is also
unchanged. In the Shapley formula \eqref{eq:shapley} every marginal
contribution of $x$ increases by exactly $a_x$, and the weights sum to one, so
$\phi_x(c')=\phi_x(c)+a_x$. Therefore
$\sigma(S;\phi(c'))=c(S)+a(S)-\phi(c)(S)-a(S)=\sigma(S;\phi(c))$.

With a common target $\alpha(S)\equiv\alpha$, the required volume $D_i$ of each
retailer does not depend on the coalition. Because every ton leaving retailer
$i$ incurs $c_i^R$, the collection term in \eqref{eq:objective} equals
$\sum_{i\in I(S)}c_i^RD_i$, which is additive. The Manufacturer's platform cost
and the per-retailer administration cost are additive by construction. This
proves Proposition~\ref{prop:ch4_strategic_equivalence}.

\subsection{The Linear Relaxation}

Write $\bar c(S)=\bar C(S)+F^G(S)$, where $\bar C(S)$ is the optimal value of
the relaxed network program (\eqref{eq:objective}--\eqref{eq:con_domain} with
$0\leq y_j,z_l\leq1$, excluding the governance term). For a platform-controlled
processor the bound $z_l\leq1$ is implied by $z_l\leq e_l(S)\leq1$ and can be
dropped; for a processor whose eligibility is the same for every coalition the
access row is a constant $z_l\leq e_l$, and when $e_l=1$ it coincides with the
bound, so one copy can be dropped. The relaxed program for
every coalition then has the same constraint matrix and differs only in its
right-hand side $b(S)$, whose entries are of three kinds:
\begin{enumerate}
\item service rows \eqref{eq:con_retailer}, with right-hand side
$D_i\mathbf{1}[i\in S]$, which adds up over retailers;
\item access rows for platform-controlled processors, with right-hand side
$\mathbf{1}[M\in S]$, contributed by the Manufacturer;
\item shared rows---flow balance, capacity-link rows, and access rows of
processors unavailable to every coalition, all with right-hand side zero; and
the bounds $y_j\leq1$ and $z_l\leq1$ for processors open to every coalition,
with right-hand side one.
\end{enumerate}
Let $u^*$ be an optimal dual solution of the grand coalition's relaxed program.
By condition~(d) of Proposition~\ref{prop:ch4_lp_balanced} a primal optimum
leaves every shared bound slack, so complementary slackness gives $u^*=0$ on
those rows; rows with right-hand side zero contribute nothing to any dual
objective. Define
\begin{equation*}
\xi_i=D_iu^*_i\quad(i\in I),
\qquad
\xi_M=\sum_{l\ \text{platform-controlled}}u^*_l .
\end{equation*}
For any coalition $S$, $\xi(S)=b(S)^{\top}u^*$. Since the dual feasible set does
not depend on $S$, $u^*$ is feasible for the dual of $S$'s program, and weak
duality gives $\xi(S)\leq\bar C(S)$. Strong duality at the grand coalition
gives $\xi(N)=\bar C(N)$. Hence $\xi$ lies in the core of $\bar C$; this is the
linear-production-game argument of \cite{Owen1975}.

For the governance term, condition~(c) writes
$F^G(S)=f(S)+F_0\mathbf{1}[I(S)\neq\emptyset]$ with member-specific terms $f$
and a setup cost $F_0\geq0$. The allocation $g_i=F_0/|I|$ for retailers and
$g_M=0$ gives $g(S)=F_0|I(S)|/|I|\leq F_0\mathbf{1}[I(S)\neq\emptyset]$ with
equality at $N$, so $g$ is in the core of the setup-cost game. Because the core
of a sum of games contains the sum of their cores, $\bar x=\xi+f+g$ lies in the
core of $\bar c$. This proves Proposition~\ref{prop:ch4_lp_balanced}.

\subsection{Indivisibility}

The relaxed program has a larger feasible set than the MILP, so
$\bar c(S)\leq c(S)$ for every coalition. The allocation $\bar x$ constructed
above satisfies $\bar x(S)\leq\bar c(S)\leq c(S)$ for every proper $S$ and
$\bar x(N)=\bar c(N)$. If $c(N)=\bar c(N)$, $\bar x$ is in the core of $c$.
Otherwise $\bar x$ is feasible in \eqref{eq:ch4_cost_of_stability} with
$\Delta=c(N)-\bar c(N)$, which bounds $\Delta^*$. This proves
Proposition~\ref{prop:ch4_indivisibility}.

\subsection{Balanced-Collection Certificate}

The least-core program minimizes $\epsilon$ subject to
$a(S)-\epsilon\leq c(S)$ for every proper $S$ (multiplier $\lambda_S\geq0$) and
$a(N)=c(N)$ (multiplier $\mu$, free). Its dual is
\begin{equation*}
\max_{\lambda\geq0,\ \mu}\ \mu\,c(N)-\sum_S\lambda_Sc(S)
\quad\text{s.t.}\quad
\sum_{S\ni x}\lambda_S=\mu\ \ \forall x\in N,
\qquad
\sum_S\lambda_S=1 .
\end{equation*}
Summing the first constraint over players gives
$\sum_S\lambda_S|S|=n\mu$, so $\mu>0$. Setting $\delta_S=\lambda_S/\mu$ yields
a balanced collection with $\sum_S\delta_S=1/\mu$, and the dual objective
becomes $\big(c(N)-\sum_S\delta_Sc(S)\big)/\sum_S\delta_S$. Conversely, every
balanced $\delta$ defines a dual-feasible pair
$\lambda=\delta/\sum_S\delta_S$, $\mu=1/\sum_S\delta_S$. The primal is feasible
and the dual is feasible, so strong duality gives \eqref{eq:ch4_certificate}.
The core is nonempty exactly when $\epsilon^*\leq0$, which proves
Proposition~\ref{prop:ch4_balanced_certificate}.

\subsection{Envelope Result for Stability}

Let $\Lambda$ be the dual feasible set above; it does not depend on the costs.
Then $\epsilon^*=E(c)=\max_{(\lambda,\mu)\in\Lambda}\mu c(N)-\sum_S\lambda_Sc(S)$
is a maximum of linear functions of the cost vector. Where the maximizer is
unique, $E$ is differentiable with $\partial E/\partial c(N)=\mu$ and
$\partial E/\partial c(S)=-\lambda_S$ (Danskin's theorem).

For each coalition, $c(S;\theta)$ is the minimum over facility configurations
of the corresponding flow-program values. Each of these values is continuous in
$\theta$, so a configuration that is uniquely optimal at $\theta$ remains
optimal nearby, and $c(S;\theta)$ coincides locally with the value of a single
linear program. With a unique dual solution, that value is differentiable and
its derivative is the shadow price of the constraint that contains $\theta$
(for a capacity $K_l$ with $z_l$ fixed at $\bar z_l$, the right-hand side is
$K_l\bar z_l$ and the derivative is the capacity shadow price times $\bar z_l$).
The chain rule gives \eqref{eq:ch4_stability_envelope}, proving
Proposition~\ref{prop:ch4_stability_envelope}. Proposition~\ref{prop:ch4_shapley_statics}
follows because \eqref{eq:shapley} is a fixed linear combination of coalition
costs: differentiating term by term gives $\partial\phi_x/\partial\theta=\phi_x(c')$,
and the slack derivative follows from \eqref{eq:ch4_core_slack}.

\subsection{Bounds on the Cost of Stability}

Program \eqref{eq:ch4_cost_of_stability} is equivalent to maximizing $a(N)$
subject to $a(S)\leq c(S)$ for every proper $S$. Its dual minimizes
$\sum_S\delta_Sc(S)$ over balanced collections, and singletons are balanced, so
both programs are feasible and
$\Delta^*=\max\{0,\ \max_{\delta}\,[c(N)-\sum_S\delta_Sc(S)]\}$. By
Proposition~\ref{prop:ch4_balanced_certificate},
$c(N)-\sum_S\delta_Sc(S)\leq\epsilon^*\sum_S\delta_S$ for every balanced
$\delta$. Because $\sum_S\delta_S|S|=n$ and $1\leq|S|\leq n-1$,
\begin{equation*}
\frac{n}{n-1}\leq\sum_S\delta_S\leq n .
\end{equation*}
If $\epsilon^*>0$, the maximizing collection $\delta^*$ of
\eqref{eq:ch4_certificate} gives
$\Delta^*\geq c(N)-\sum_S\delta^*_Sc(S)=\epsilon^*\sum_S\delta^*_S\geq
\tfrac{n}{n-1}\epsilon^*$, and every $\delta$ gives
$c(N)-\sum_S\delta_Sc(S)\leq n\epsilon^*$, so $\Delta^*\leq n\epsilon^*$. When
$\delta^*$ also solves the dual of \eqref{eq:ch4_cost_of_stability}, the first
inequality holds with equality. If $\epsilon^*\leq0$ every gap is nonpositive
and $\Delta^*=0$. The integrality-gap bound is
Proposition~\ref{prop:ch4_indivisibility}. This proves
Proposition~\ref{prop:ch4_cos_bounds}.

\section{Computational Method}

The reproducible implementation is in \texttt{chapter4\_model/}. It uses
Python~3.14 and the OR-Tools SCIP interface, with CBC as a fallback. The command
is
\begin{verbatim}
python solve_chapter4.py --data case_data.json `
    --output generated --solver SCIP
\end{verbatim}

For each scenario the program performs the following steps:
\begin{enumerate}
\item Enumerate all $2^6=64$ coalitions of
$N=\{M,R_1,R_2,R_3,R_4,R_5\}$.
\item Build and solve the coalition-specific MILP with service, access, flow,
and capacity constraints.
\item Recompute the objective independently from the returned variable values
and require agreement within $10^{-4}$ thousand dollars.
\item Compute exact Shapley allocations from all coalition costs and verify
efficiency and the savings-game identity.
\item Evaluate all 62 nonempty proper-coalition core constraints.
\item Solve the least-core LP, followed by the L1 projection in
Equation~\eqref{eq:ch4_constrained_projection}.
\item For the structural diagnosis, re-solve all coalitions with the binary
facility decisions relaxed, solve the dual of the least-core LP to recover the
balanced-collection certificate (and test that it is unique), and solve the
Cost-of-Stability LP \eqref{eq:ch4_cost_of_stability}.
\item For the comparative statics, fix each coalition's optimal facility
configuration, re-solve the remaining flow LP to obtain shadow prices, and
compare the prediction of Proposition~\ref{prop:ch4_stability_envelope} with a
central finite difference of the fully re-solved game.
\item Repeat the diagnosis over the high-standard processor capacity, the
formal-service target, and a mean-preserving volume dispersion.
\item Stop with an error if any exported result contradicts
Propositions~\ref{prop:ch4_strategic_equivalence}--\ref{prop:ch4_cos_bounds}:
a relaxed game with an empty core where condition~(d) holds, a Cost of
Stability outside its bounds, a change in stability when additive costs are
removed, or an envelope prediction that disagrees with the finite difference.
\item Export CSV, JSON, LaTeX table fragments, pgfplots data, and a SHA-256
manifest.
\end{enumerate}
The full run, including the diagnosis, takes about two minutes.

Variable costs enter the objective in dollars per ton and are divided by 1,000;
fixed and reported costs are in thousands of U.S. dollars. The dataset is a
transparent numerical illustration, not an empirical estimate.

\section{Complete Stylized Input Data}

\begin{table}[htbp]
\centering
\caption{Candidate Consolidation-Center Inputs}
\label{tab:app_ch4_centers}
\begin{tabular}{lccc}
\toprule
Center & Capacity (tons) & Fixed cost (thousand USD) & Handling (USD/ton) \\
\midrule
% Generated parameter rows; do not edit.
NV & 6,500 & 300 & 35 \\
TN & 6,000 & 260 & 32 \\
MI & 5,000 & 230 & 30 \\
GA & 5,000 & 220 & 28 \\
\bottomrule

\end{tabular}
\end{table}

\begin{table}[htbp]
\centering
\caption{Processor Inputs}
\label{tab:app_ch4_processors}
\resizebox{\textwidth}{!}{%
\begin{tabular}{lrrrrrr}
\toprule
Processor & Capacity & Fixed cost & Processing & Recovery efficiency & Base recovery value & Access rule \\
\midrule
High-standard $H$ & 10,000 & 550 & 410 & 0.90 & 330 & Requires $M\in S$ \\
Standard $L$ & 15,000 & 280 & 300 & 0.65 & 210 & Available to retailer coalitions \\
\bottomrule
\end{tabular}%
}
\end{table}

\begin{table}[htbp]
\centering
\caption{Retailer-to-Center Transportation Cost (USD/ton)}
\label{tab:app_ch4_retailer_transport}
\begin{tabular}{lrrrr}
\toprule
 & NV & TN & MI & GA \\
\midrule
$R_1$ & 45 & 155 & 135 & 170 \\
$R_2$ & 85 & 55 & 90 & 75 \\
$R_3$ & 125 & 75 & 25 & 70 \\
$R_4$ & 150 & 55 & 80 & 20 \\
$R_5$ & 170 & 105 & 60 & 95 \\
\bottomrule
\end{tabular}
\end{table}

\begin{table}[htbp]
\centering
\caption{Center-to-Processor Transportation Cost (USD/ton)}
\label{tab:app_ch4_processor_transport}
\begin{tabular}{lrr}
\toprule
 & High-standard $H$ & Standard $L$ \\
\midrule
NV & 40 & 55 \\
TN & 55 & 30 \\
MI & 25 & 50 \\
GA & 60 & 25 \\
\bottomrule
\end{tabular}
\end{table}

The baseline formal-service target is $\alpha(S)=1$ for every coalition with
retail return volume, so coalitions are compared under the same service
obligation. The standard processor is available to all such coalitions; the
high-standard processor requires the Manufacturer's platform. The Manufacturer
platform costs 500 thousand dollars. For the Manufacturer singleton, this is
interpreted as the annual outside-option cost of maintaining platform readiness
and the processor relationship without physical retailer volume. A coalition
with retailer volume also incurs 50 thousand dollars of setup cost plus 25
thousand dollars per participating retailer. Processor prices and recovered
values are fixed by processor type rather than by an abstract capability score.
Recovery efficiency is used to report recovered output and determine the
recovered-material credit; the baseline does not impose a minimum recovery-yield
constraint. The JSON file is the authoritative machine-readable source.
